\section{Setting and main theorem}
Fix a probability space $(\Omega,\F,\Pp)$ and a filtration $(\F_t)_{t\geq0}$
satisfying the usual conditions. The price $S$ is a bounded real-valued
adapted c\`adl\`ag semimartingale. Boundedness means that one deterministic
constant $B$ satisfies $\sup_{t\geq0}|S_t|\leq B$ almost surely.
The Lean source chooses an everywhere c\`adl\`ag representative; the bound is
required almost surely. The public theorem also explicitly includes
\path{SigmaFiniteFiltration}, which holds because restrictions of a probability
measure to the filtration are finite.

Semimartingale behavior is expressed by the good-integrator property.
An elementary predictable strategy $J$ has a pathwise finite-sum integral;
its value at each finite time is required to be measurable up to completion.
If a sequence $J^n$ satisfies
\[
 \forall\varepsilon>0\ \exists N\ \forall n\geq N\ \forall(t,\omega),
 \qquad |J^n_t(\omega)|\leq\varepsilon,
\]
then $(J^n\cdot S)_T\to0$ in probability for every finite $T$.
Section~\ref{sec:graph} constructs integrals of general predictable coefficients
from these assumptions. We write $H\cdot S$ for such an integral, always with
initial value zero.

\begin{definition}
An integrable strategy $H$ is $a$-admissible if, on one set of probability one,
\[
 (H\cdot S)_t\geq-a\qquad\text{for every }t\geq0.
\]
Let $K_0$ be the terminal gains of strategies admissible for some finite
constant and possessing a finite almost-sure limit as $t\to\infty$.
Let $K_1$ be the terminal gains of such $1$-admissible strategies. Set
\[
 C_0=K_0-L^0_+,\qquad C=C_0\cap L^\infty.
\]
\end{definition}
The set $K_0$ is a convex cone and $K_1$ is convex, with $K_1\subseteq K_0$.
The cone $C_0$ is downward solid for the almost-sure order. NFLVR means
\[
 \overline C^{\,\|\cdot\|_\infty}\cap L^\infty_+=\{0\}.
\]
A set $D\subseteq L^0$ is Fatou closed if $f_n\in D$, a common constant lower
bound on $f_n$, and $f_n\to f$ almost surely imply $f\in D$.

For the implementation, distinguish $C_0\subseteq L^0$ from its raw-function
extension
\[
 \widehat C_0=\{f:\Omega\to\mathbb R:\exists g\in K_0,\ f\leq g\text{ a.s.}\}.
\]
Here $K_0$ denotes raw representatives of terminal gains.
Lean uses $\widehat C_0$. If $f$ is measurable up to completion, then so is
$g-f$, because terminal gains are measurable up to completion. Thus
$f\in\widehat C_0$ if and only if $f=g-h$ almost surely for some $g\in K_0$
and an $h$ measurable with respect to completion and nonnegative almost surely.
Passing to almost-sure
classes recovers exactly $C_0=K_0-L^0_+$. The raw Fatou-closedness result
therefore implies the stated $L^0$ result, and its intersection with $L^\infty$
gives exactly the same cone $C$ and NFLVR condition.

\begin{theorem}[Delbaen--Schachermayer]\label{thm:main}
Assume NFLVR. Then $C_0$ is Fatou closed, and $C$ is closed in the topology
$\sigma(L^\infty,L^1)$.
\end{theorem}
The integral in this statement is the general integral graph of
Section~\ref{sec:graph}: bounded coefficients are integrated in a local
$M^2\oplus A^1$ completion, and general gains arise as limits of coefficient
truncations uniformly over elementary tests. A terminal claim is the
continuous-time almost-sure limit of its own gain process.
Lean's definition of Fatou closedness normalizes the common lower bound to
$-1$. For a cone stable under positive scaling, divide a sequence bounded below
by $-b$ by $\max(1,b)$ and then scale its limit back.

This is the real-valued version of \cite[Theorem~4.2]{DS}.
The essential probabilistic step realizes an almost-sure maximal element of
$\clprob{K_1}$ as an admissible integral's terminal value.
Weak-star closedness then follows by functional analysis.

\subsection{Proof structure}
First construct the general integral and its closedness under regular Émery
limits. This construction does not assume market closedness.
NFLVR gives boundedness in probability and forward convex compactness of $K_1$.
For a maximal closure claim $h$, chronological pasting yields integral gains
converging uniformly over all time to a process $X$ with $X_\infty=h$.
Choose one equivalent measure $Q_*$ and decompose each selected gain separately.
Common convex weights make the martingale components Cauchy by tail estimates
and Hilbert convexification. Hahn improvements and maximality make the
finite-variation components Cauchy. These estimates imply Émery convergence of
the same gains. Integral closedness yields $h\in K_1$, from which the two
closedness conclusions follow.
The market argument ends in Section~\ref{sec:fatou}.
Its boundary with the appendices consists of the seven analytic propositions (twenty clauses) and two data definitions in Section~\ref{sec:graph}.
Each clause corresponds to one input in the Lean isolation check. The main proof uses only these conclusions.
A reader wishing to check an input can turn to its indicated appendix proof.

For decomposition, the appendix explicitly cites the bounded
Bichteler--Dellacherie characterization and existence of dual predictable
projections as standard results, with their hypotheses specified.
Lean constructs these too. Integral-range closedness, simultaneous
realization by one coefficient, and original-market component estimates
are proved within this account. The final appendix relates the mathematical
stages to Lean declarations.

\subsection{Correspondence with the original proof}
\label{sec:original-proof}
The main argument follows Delbaen--Schachermayer\cite[Section 4]{DS}:
selection of a maximal closure claim, pasting, common convexification under
an equivalent measure, and a maximality contradiction using Hahn improvements.
The table matches the content used in the main proof; it does not restate
every original result in its full generality.

\begin{center}
\small
\begin{tabular}{L{0.24\linewidth}L{0.66\linewidth}}
\toprule
Original paper & Corresponding argument here \\
\midrule
Lemma A1.1 & Proposition~\ref{prop:forward}: forward convex compactness in a lower-bounded convex set \\
Lemma 4.3 & Proposition~\ref{prop:maximal}: existence of a maximal terminal claim \\
Lemma 4.5 & Proposition~\ref{prop:all-time-cauchy}: all-time Cauchy estimates from maximality \\
Lemma 4.7 & Proposition~\ref{prop:indexed-maximal}: class bounds for martingale maxima \\
Lemma 4.8 & Proposition~\ref{prop:convex-tails} and its corollary: convex stopped-tail estimates \\
Lemma 4.9 & Proposition~\ref{prop:test-uniform-tails}: test-uniform tail estimates \\
Lemma 4.10 & Proposition~\ref{prop:martingale-convexification}: common martingale convexification \\
Lemma 4.11 & Section~\ref{sec:fv-cauchy}: finite-horizon FV variation Cauchy estimates from Hahn improvements \\
Role of Lemma A3.1 & Lemma~\ref{lem:persistent-improvement}: persistent terminal improvement contradicts maximality \\
Final use of M\'emin's theorem & Proposition~\ref{input:realization} and Section~\ref{sec:realization}: realizing the limit as an original-price integral \\
Use of Theorem 2.1 & Proposition~\ref{prop:weakstar}: from Fatou to weak-star closedness \\
\bottomrule
\end{tabular}
\end{center}

We obtain maximal elements by convex compactness and maximization of a
functional, replacing the original transfinite induction. Component estimates
allow separate gain decompositions, but every contradiction strategy returns
to the same original market. For martingale components we first establish
Cauchy estimates and identify the limit later; for FV components we prove the
needed finite-horizon estimates and combine them with the separately retained
infinite-time terminal limit. Proposition~\ref{input:finite} supplies the finite-scale analytic estimates needed in Lemmas 4.7--4.9.
The integral-range closedness cited by the original paper enters the proof of terminal realization, Proposition~\ref{input:realization}. These are reorganizations
preserving the roles of the original argument, not a claim to a new proof
strategy for the closedness theorem.
